% ===== spotify-frames.tex: the Spotify case in three frames, asked for by Carlos on 2026-11-06 =====
% The frames of the case and their talk track. Since 2026-11-09 two decks input this file, each as one section:
% block2-matrix-factorization.tex, the deck of its Friday block, and weekend2-cases.tex, the weekend deck.
% His dictated text: (1) little data at the launch of audiobooks, the purchase before streaming, a quarter of the
% users with three quarters of the hours, over 70% of the first listeners already on podcasts; (2) 128 numbers per
% listener and per audiobook, the shelf "Audiobook for you" by inner product, two networks that compute the numbers
% every day from inputs on both sides and a graph of podcasts and audiobooks; (3) the A/B test of 11.5 million users
% in three groups and the two gains in new audiobook starts. Diagrams on frames 2 and 3, as he asked.
% Simplified the same day after his "This is too complicated for me ... Make them 25% simpler": frame 1 lost its
% bars (the text says the same), frame 2 has a graph of four items and shorter labels, frame 3 has no reference
% line. README.md, section "One deck per case", has the counts.
% No number is typed here: numbers.tex is written by tools/cases_numbers.py (block CASE 2), which keeps the
% sentence of each source. A "Say" comment above each frame is the talk track. Build and check from this folder:
%   python3 tools/cases_numbers.py
%   ../style/tools/build.sh block2-matrix-factorization.tex
%   python3 tools/check_numbers.py block2-matrix-factorization.pdf tools/expected-block2-matrix-factorization.txt
\newcommand{\spsource}{De Nadai and others (Spotify), WWW 2024 Companion.}

% ===== Frame 1. Little data to learn from =====
% Say: When Spotify launched audiobooks it had little data to learn from. Each audiobook had to be bought before
% it could be streamed, so few people tried one, and the listening was concentrated: about a quarter of the
% audiobook listeners made three quarters of the listening hours. [click] But Spotify knew something else about
% its users: over 70 percent of the first audiobook listeners already listened to podcasts. Podcast tastes were
% the way in.
\begin{frame}{Spotify's audiobooks: little data to learn from}
  \begin{tikzpicture}[x=1mm, y=1mm]
    \path[use as bounding box] (0,1) rectangle (130,-54);
    \node[body, anchor=north west, text width=64mm] at (0,0)
      {When Spotify launched audiobooks it had little data to learn from.\par\vspace{2.4mm}
       Each audiobook had to be bought before it could be streamed, and about a quarter of the audiobook listeners
       made three quarters of the listening hours.};
    \node[big, anchor=north west, visible on=<2->] at (76,-1) {Over \spFromPodcasts\%};
    \node[body, anchor=north west, text width=54mm, visible on=<2->] at (76,-10)
      {of the first audiobook listeners already listened to podcasts.};
    \node[src, anchor=north west, text width=130mm] at (0,-50.6) {\spsource};
  \end{tikzpicture}
\end{frame}

% ===== Frame 2. Two networks, a graph and an inner product =====
% Say: Every listener and every audiobook gets a list of 128 numbers. The shelf "Audiobook for you" on the home page
% picks from the audiobooks whose numbers have the highest inner product with the listener's. In the matchmaker we
% learned the numbers directly. Here two networks, trained again every day, compute them from inputs on both
% sides: for the listener, what they listened to, audiobooks, podcasts and artists, and their age and country; for
% the audiobook, its language, genre, title and description. The audiobooks' numbers are computed once a day, the
% listener's at each visit. [click] And a graph that links two items, a podcast and an audiobook or two of a kind,
% whenever one user streamed both. It feeds both networks, so podcast tastes carry over to audiobooks.
\begin{frame}{\spDims{} numbers for every listener and every audiobook}
  \begin{tikzpicture}[x=1mm, y=1mm]
    \path[use as bounding box] (0,1) rectangle (130,-54);
    \node[body, anchor=north west, text width=130mm] at (0,0)
      {The matchmaker learned the numbers. Here two networks compute them.};
    % inputs
    \node[anchor=north, inner sep=0, font=\small, align=center] at (22,-9) {Listening history,\\ age, country};
    \node[anchor=north, inner sep=0, font=\small, align=center] at (110,-9) {Language, genre,\\ title, description};
    % the graph (click 2): two podcasts (circles), two audiobooks (squares); links of both kinds
    \node[lbl, anchor=south, visible on=<2->] at (57,-10.6) {podcasts};
    \node[lbl, anchor=south, visible on=<2->] at (75,-10.6) {audiobooks};
    \foreach \b/\d in {13.5/13.5, 19.5/13.5, 19.5/19.5}
      {\draw[draw=nib, line width=0.5pt, visible on=<2->] (57,-\b) to (75,-\d);}
    \draw[draw=nib, line width=0.5pt, visible on=<2->] (55.65,-13.5) to[bend right=75] (55.65,-19.5);
    \draw[draw=nib, line width=0.5pt, visible on=<2->] (76.35,-13.5) to[bend left=75] (76.35,-19.5);
    \foreach \ny in {13.5,19.5} {
      \draw[draw=nib, line width=0.6pt, fill=paper, visible on=<2->] (57,-\ny) circle (1.35);
      \fill[hatch fine, visible on=<2->] (73.65,-\ny-1.35) rectangle (76.35,-\ny+1.35);
      \draw[draw=brand, line width=0.6pt, line join=miter, visible on=<2->] (73.65,-\ny-1.35) rectangle (76.35,-\ny+1.35);}
    \node[elab dim, anchor=north, visible on=<2->] at (66,-22) {linked by a shared listener};
    % the networks
    \node[net] (lnet) at (22,-29.5) {Listener network};
    \node[net] (anet) at (110,-29.5) {Audiobook network};
    \draw[flow] (22,-17.4) to (lnet.north);
    \draw[flow] (110,-17.4) to (anet.north);
    \draw[flow, visible on=<2->] (66,-25.6) to (66,-29.5) to (lnet.east);
    \draw[flow, visible on=<2->] (66,-29.5) to (anet.west);
    % the numbers, the inner product and the shelf
    \draw[flow] (lnet.south) to (22,-35.8);
    \draw[flow] (anet.south) to (110,-35.8);
    \hatchcells{1-}{15.1}{-38}{4}
    \hatchcells{1-}{103.1}{-38}{4}
    \node[elab, inner sep=1pt] (ip) at (66,-38) {inner product};
    \draw[flow] (31.4,-38) to (ip.west);
    \draw[flow] (100.6,-38) to (ip.east);
    \node[chip, font=\normalsize, minimum height=5.6mm, minimum width=36mm, inner sep=1.5pt] (shelf) at (66,-45.6)
      {Audiobook for you};
    \draw[flow] (ip.south) to (shelf.north);
    \node[src, anchor=north west, text width=130mm] at (0,-50.6) {\spsource};
  \end{tikzpicture}
\end{frame}

% ===== Frame 3. The A/B test =====
% Say: The value. Spotify ran an A/B test on 11.5 million monthly users, split at random into three groups: the
% model already in production, two networks without the graph, and the full model with the graph.
% [click] The measure: the share of listeners who started a new audiobook. Against the production model it rose by
% 46.8 percent with the full model and by 23.9 percent with the networks alone: the graph roughly doubled the
% gain. Both are relative changes. By the time of the paper, 2024, the full model was in production for every
% eligible user.
\begin{frame}{The value: an A/B test}
  \begin{tikzpicture}[x=1mm, y=1mm]
    \path[use as bounding box] (0,1) rectangle (130,-54);
    \node[chip, font=\normalsize, text width=25mm, align=center, minimum width=28mm, inner sep=2pt] (users) at (14.5,-28)
      {\spUsers{} million monthly users};
    \node[lbl, anchor=north] at (14.5,-34) {split at random};
    \foreach \gy/\gname/\gtext in {-13/g1/{Production model}, -28/g2/{Two networks},
                                   -43/g3/{Two networks and the graph}}
      {\node[chip, font=\small, text width=33mm, align=center, minimum width=36mm, minimum height=7.6mm,
         inner sep=1.5pt] (\gname) at (52,\gy) {\gtext};}
    \draw[draw=nib, line width=0.7pt, line join=miter] (users.east) to (30,-28);
    \draw[draw=nib, line width=0.7pt, line join=miter] (30,-13) to (30,-43);
    \foreach \gy in {-13,-28,-43} {\draw[flow] (30,\gy) to (33.9,\gy);}
    % the measure and the bars (click 2): the production model's bar is 22 mm long
    \node[measure, anchor=north west, text width=56mm, visible on=<2->] at (74,0)
      {Listeners starting a new audiobook, relative change};
    \draw[bar plain, visible on=<2->] (74,-15.5) rectangle (96,-10.5);
    \fill[hatch, visible on=<2->] (74,-30.5) rectangle ({74+22*\spIdxPlain/100},-25.5);
    \draw[brand, line width=0.7pt, line join=miter, visible on=<2->] (74,-30.5) rectangle ({74+22*\spIdxPlain/100},-25.5);
    \node[anchor=west, inner sep=0, font=\large, visible on=<2->] at ({76.5+22*\spIdxPlain/100},-28) {+\spStartPlain\%};
    \fill[hatch, visible on=<2->] (74,-45.5) rectangle ({74+22*\spIdxFull/100},-40.5);
    \draw[brand, line width=0.7pt, line join=miter, visible on=<2->] (74,-45.5) rectangle ({74+22*\spIdxFull/100},-40.5);
    \node[big, anchor=west, visible on=<2->] at ({76.5+22*\spIdxFull/100},-43) {+\spStart\%};
    \node[src, anchor=north west, text width=130mm] at (0,-50.6) {\spsource};
  \end{tikzpicture}
\end{frame}

